\newpage

\section{Exercices}
 {
  \small
  \begin{multicols}{2}
	  \begin{exo}[Analyse d'affectations]
		  Soient les instructions suivantes exécutées successivement :

		  \begin{lstlisting}
a = 3
b = 7
a = a + b
b = a - b
a = a - b
\end{lstlisting}

		  Quelles sont les valeurs finales des variables \lstinline|a| et \lstinline|b| ?
	  \end{exo}

	  \begin{exo}[Calculs arithmétiques]
		  Soient \lstinline|x = 17| et \lstinline|y = 5|. Donner le résultat et le type des expressions suivantes :
		  \begin{enumerate}[label=\alph*)]
			  \item \lstinline|x / y|
			  \item \lstinline|x // y|
			  \item \lstinline|x % y|
			  \item \lstinline|x ** 2|
		  \end{enumerate}
	  \end{exo}

	  \begin{exo}{}
		  Expliquer de façon claire et concise ce que fait le programme suivant.

		  \begin{lstlisting}
a=int(input("Entrez un entier"))
b=a
a=a*b
print(a)
\end{lstlisting}
	  \end{exo}

	  \begin{exo}{}
		  Qu'affichera le programme suivant ?

		  \begin{lstlisting}
a=1
b=2
c=a+b
a=b+c
b=c+a
print(b)
\end{lstlisting}
	  \end{exo}

	  \begin{exo}{}
		  Que fait le programme suivant ?

		  \begin{lstlisting}
a=int(input("Entrez un entier"))
b=a+a
a=a+b
print(a)
\end{lstlisting}
	  \end{exo}

	  \begin{exo}{}
		  Qu'affichera le programme suivant ?

		  \begin{lstlisting}
a=3
for i in range(1,1001):
    a=a+0.5
print(a)
\end{lstlisting}
	  \end{exo}

	  \vfill~\columnbreak

	  \begin{exo}{}
		  Que fait le programme suivant ?

		  \begin{lstlisting}
b=1
a=int(input("Entrez un nombre"))
for i in range(1,4):
    b=b*a
print(b)
\end{lstlisting}
	  \end{exo}

	  \begin{exo}{}
		  Compléter le programme suivant pour qu'il calcule la somme des inverses des entiers de $1$ à $1000$ :

		  \[
			  1 + \frac12 + \frac13 + \frac14 + \cdots + \frac1{1000}
		  \]

		  \begin{lstlisting}
somme=0
for i in range(1,1001):
    somme=somme+.........
print(somme)
\end{lstlisting}
	  \end{exo}

	  \begin{exo}{}
		  Écrivez un programme qui demande à l'utilisateur d'entrer un entier naturel non nul $n$ et qui affiche la factorielle de $n$.

		  La factorielle d'un entier $n$, notée $n!$, est le produit de tous les entiers de $1$ à $n$ inclus.

		  \[
			  5! = 1 \times 2 \times 3 \times 4 \times 5 = 120.
		  \]
	  \end{exo}

	  \begin{exo}{}
		  Que fait le programme suivant ?
		  {\footnotesize
		  \begin{lstlisting}
s=0
n=int(input("Entrez un entier > 0"))
for i in range(1,n+1):
    s=s+i
print(s)
\end{lstlisting}
		  }
	  \end{exo}

	  \begin{exo}{}
		  Que fait le programme suivant ?
		  {\footnotesize
		  \begin{lstlisting}
s=0
n=int(input("Entrez un entier > 0"))
p=int(input("Entrez-en un deuxième"))
for i in range(1,n+1):
    for j in range(1,p+1):
        s=s+1
print(s)
\end{lstlisting}
		  }
	  \end{exo}

	  \begin{exo}[Structure conditionnelle]
		  Écrire une fonction Python \lstinline|maximum(a, b)| qui prend deux nombres \lstinline|a| et \lstinline|b| en paramètres et renvoie le plus grand des deux sans utiliser la fonction native \lstinline|max()|.
	  \end{exo}

	  \vfill~\columnbreak

	  \begin{exo}{}
		  Le programme suivant tire au sort $20$ entiers entre $1$ et $100$ et affiche leur somme.

		  Modifiez-le pour qu'il affiche non plus la somme, mais le plus grand d'entre eux.

		  \begin{lstlisting}
from random import *

s=0
for i in range(1,21):
    a=randint(1,100)
    s=s+a
print(s)
\end{lstlisting}
	  \end{exo}

	  \begin{exo}{}
		  Modifiez le programme de l'exercice précédent pour qu'il affiche à la fin la moyenne des entiers tirés au sort.
	  \end{exo}

	  \begin{exo}[Boucle \texttt{while} et seuil]
		  On considère la suite de nombres définie par son premier terme $U_0 = 100$ et la relation $U_{n+1} = 0{,}8 U_n + 5$.

		  Écrire un programme Python déterminant le plus petit entier $n$ tel que $U_n < 30$.
	  \end{exo}

	  \begin{exo}{}
		  Qu'affichera le programme suivant ?

		  \begin{lstlisting}
a=30
n=1
while n<=a:
    n=2*n
print(n)
\end{lstlisting}
	  \end{exo}

	  \begin{exo}{}
		  Que fait le programme suivant ?

		  \begin{lstlisting}
a=input("Entrez un entier > 0")
a=int(a)
n=1
while n<=a:
    n=2*n
print(n)
\end{lstlisting}
	  \end{exo}

	  \begin{exo}{}
		  Un puits contient une hauteur d'eau de $10$ m. Chaque jour, son niveau diminue de $2$\%.

		  Compléter l'algorithme suivant pour qu'il calcule le nombre de jours qu'il faudra pour que le niveau descende en dessous de $1$ m.

		  \begin{lstlisting}
hauteur=10
nb_jours=0

while hauteur>1:
    nb_jours=............
    hauteur=............

print(nb_jours)
\end{lstlisting}
	  \end{exo}

	  \vfill~\columnbreak

	  \begin{exo}[Milieu et distance dans le plan]
		  On considère deux points $A(x_A, y_A)$ et $B(x_B, y_B)$ dans un repère orthonormé.

		  Écrire une fonction Python \lstinline|milieu(xA, yA, xB, yB)| qui renvoie les coordonnées $(x_I, y_I)$ du milieu $I$ du segment $[AB]$ sous forme d'un couple de valeurs.
	  \end{exo}

	  \begin{exo}{}
		  Qu'affichera le programme suivant ?

		  \begin{lstlisting}
def f(x):
    y=x**2+x+1
    return y

y=f(5)
print(y)
\end{lstlisting}
	  \end{exo}

	  \begin{exo}{}
		  Que fait la fonction suivante ?

		  \begin{lstlisting}
def ultime(a,b,c):
    if b>a:
        a=b
    if c>a:
        a=c
    return a
\end{lstlisting}
	  \end{exo}

	  \begin{exo}{}
		  Une urne contient $3$ boules blanches et $3$ boules noires. On effectue deux tirages avec remise.

		  Écrire une fonction qui simule la répétition à $n$ reprises de cette expérience et affiche le nombre de cas où le joueur a gagné (deux boules blanches).
	  \end{exo}

	  \begin{exo}{}
		  Dans le cadre d'un jeu de hasard, un joueur lance un dé.

		  S'il obtient \epsdice{5} ou $\epsdice{6}$, il gagne, sinon il perd.

		  La fonction suivante simule $n$ fois l'expérience et renvoie le nombre de cas où le joueur a gagné.

		  Modifier cette fonction afin qu'elle renvoie la \textbf{fréquence} des victoires.

		  Que se passe-t-il lorsque $n$ devient très grand ?
		  {\footnotesize
		  \begin{lstlisting}
from random import *

def echantillon(n):
    gagne=0
    for i in range(1,n+1):
        a=randint(1,6)
        if a>=5:
            gagne=gagne+1
    return gagne
\end{lstlisting}
		  }
	  \end{exo}
	  \vfill~

  \end{multicols}
 }
